\documentclass[sigconf]{acmart}
\settopmatter{printacmref=false, printccs=false, printfolios=true}
\renewcommand\footnotetextcopyrightpermission[1]{}
\usepackage{booktabs}
\usepackage{xcolor}
\newcommand{\todo}[1]{\textcolor{red}{[TODO: #1]}}

\begin{document}
\title{Performance Engineering of a PySpark Pipeline on NYC Yellow Taxi Trips}

\author{Tomás Neves}
\email{\todo{email Tomás}}
\author{Hugo Feteira}
\email{h.feteira@campus.fct.unl.pt}

\begin{abstract}
We build a PySpark pipeline over the NYC TLC Yellow Taxi trips (2019, 84.6\,M records) with three mandatory
operations (cleaning, temporal aggregation, zone ranking) and two optional ones (origin--destination flows and
demand-forecast features). Starting from a 1398\,s baseline on a 3-worker Docker cluster, a profile-driven loop
found that the TLC files hold a single Parquet row group each, which caps read parallelism at one task per file.
We also found that every operation re-reads and re-cleans the raw data. Materialising the cleaned trips with small
row groups, aggregating before lookup joins, and caching only small reused intermediates bring the pipeline to
\todo{final}\,s. Caching the full trip table instead crashed the executors with out-of-memory errors.
\end{abstract}

\keywords{Apache Spark, performance engineering, data skew, Parquet, caching}
\maketitle

%--------------------------------------------------------------------
\section{Introduction}
\textbf{Dataset.} NYC TLC Yellow Taxi trip records, one Parquet file per month. We measure on 2019 (12 files,
84.6\,M trips, \todo{GB}). \textbf{Tasks.} M1 cleaning with per-criterion removal counts; M2 trips and mean
fare/duration/distance per (borough, weekday, hour); M3 monthly zone ranking with a 3-month rolling share and
volume-anomaly flags; O2 origin--destination flows, the top-20 routes per borough, and routes whose share
reverses direction between Q1 and Q3; O3 demand features per (zone, weekday, hour): historical mean and variance,
4-week lag, and a demand index.
\textbf{Why it is interesting.} The workload mixes wide aggregations, window functions, self-joins and lookup
joins over strongly skewed keys (Manhattan $\approx$60\,\% of trips). Several operations reuse the same cleaned
input, so data layout and reuse dominate run time, not arithmetic.

%--------------------------------------------------------------------
\section{Pipeline}
\texttt{src/main.py} runs M1, then each operation as a module with the same interface
\texttt{run(trips, zones, args, opt)}. Every operation runs inside the profiler of Lab~3. Its plan
(\texttt{explain("extended")}) and metrics go to \texttt{results/}. Optimisations are switched on with
\texttt{--opt}, so every variant is comparable with the baseline in the same code base.
\textbf{M1} reads each file separately and casts it to a fixed schema (\texttt{passenger\_count} and
\texttt{RatecodeID} change type in mid-2018). It computes six removal rules in \emph{one} aggregation (one read,
not six \texttt{filter().count()}). It removes 3.80\,M trips (4.5\,\%).
\textbf{M3} builds a full zone$\times$month grid (\texttt{crossJoin} + left join) so empty months exist before
the \texttt{rangeBetween} windows are applied. The anomaly flag compares each month with the three \emph{previous}
months (with the current month inside, $|z|$ cannot exceed $2/\sqrt{3}\approx1.15$) and needs a rolling mean
$\ge$30 trips. \textbf{O3} computes mean and variance from $\sum x$ and $\sum x^2$,
dividing by the calendar count of each weekday, so hours without trips count as zero. The 4-week lag is a
self-join on \texttt{date\_add(28)}; the demand flag needs a mean $\ge$1 trip per hour.
\textbf{Decisions.} We used the DataFrame API throughout: it composes with Python helpers (the rule dictionary,
the \texttt{--opt} switches) and produces the same Catalyst plans as SQL.
\todo{Figure: simplified physical plan of O2 (Exchange / BroadcastHashJoin / Window).}

%--------------------------------------------------------------------
\section{Profiling and Baseline}
\textbf{Setup.} MacBook Pro, Intel i7-6920HQ (4C/8T), 16\,GB RAM. Docker VM: 6 vCPU, 8\,GB. Spark 4.2.0 standalone:
1 master + 3 workers, 2 cores and 1500\,MB each (executor 1200\,MB, driver 1\,GB). Default Spark configuration
(AQE on, 200 shuffle partitions). Data: 2019 full year. Median of 3 runs.
\textbf{Noise.} The VM is shared with unrelated containers, so identical runs varied up to $\pm$40\,\%
(baseline: 970, 1398, 1871\,s). For per-operation comparisons we therefore also report shuffle bytes, input
records and task counts, which do not depend on the machine's load.

\begin{table}[h]
\caption{Baseline, 2019, median of 3 runs.}
\label{tab:baseline}
\begin{tabular}{lrrrrrr}
\toprule
 & M1 & M2 & M3 & O2 & O3 & total \\
\midrule
wall time (s) & 86 & 161 & 101 & 337 & 295 & 1398 \\
tasks & 157 & 146 & 198 & 840 & 721 & \\
input records (M) & 84.6 & 84.6 & 84.6 & 197.6 & 169.2 & \\
\bottomrule
\end{tabular}
\end{table}

\textbf{Bottlenecks found.} (1)~\emph{Skewed reads}: in the M3 read stage, the slowest task took 75\,s against a
2.7\,s median (28$\times$). \texttt{spark\_partition\_id()} showed that all rows of a file land in one partition.
Every TLC file has a \emph{single} row group, so of the 3--6 splits Spark makes per file only one reads data, and
read parallelism equals the number of files. (2)~\emph{Repeated work}: no operation reuses another's result, so
each one re-reads and re-cleans the raw files. The plans contain $\approx$10 \texttt{FileScan}s per run; input
records for O2 and O3 are 2--2.3$\times$ the dataset. (3)~\emph{Wide joins in O2}: two
\texttt{BroadcastHashJoin}s to the zone lookup over every trip, followed by a shuffle carrying four string
columns.

%--------------------------------------------------------------------
\section{Optimisation Study}

\subsection{OPT-1: materialise cleaned trips with small row groups}
\emph{Bottleneck:} skewed single-row-group reads, repeated cleaning. \emph{Hypothesis:} writing the cleaned
trips once with \texttt{parquet.block.size}=32\,MB gives several tasks per file and removes the re-reads.
\emph{Change:} \texttt{--opt materialize} in M1. \emph{Result:} total 1398$\to$661\,s ($-$53\,\%). M1 pays
+217\,s for the write. M2, O2 and O3 drop by 61--70\,\% (O2 tasks 840$\to$242). Outputs are identical to the
baseline (\texttt{exceptAll} both ways; M2 differs by $10^{-14}$, from floating-point summation order).

\subsection{OPT-2: aggregate before the lookup joins (O2)}
\emph{Bottleneck:} joins and shuffle over 80\,M wide rows. \emph{Hypothesis:} counting by
(PU, DO, year, month) first leaves $\le$263$^2$ routes per month to join; Catalyst does not push a
\texttt{groupBy} below a join. \emph{Change:} \texttt{--opt agg\_first}. \emph{Result} (on top of OPT-1):
O2 125$\to$63\,s, executor CPU 360$\to$132\,s, shuffle write 74.8$\to$39.0\,MB; total 661$\to$567\,s.

\subsection{OPT-3: cache the cleaned trips (rejected)}
\emph{Hypothesis:} \texttt{persist(MEMORY\_AND\_DISK)} removes the re-reads without OPT-1's write.
\emph{Result:} M2, which fills the cache, took 789\,s (baseline 161\,s). M3 then lost 8 executors to
\texttt{OutOfMemoryError} and the job aborted. Storage and execution share the 1200\,MB heap, and the decoded
trips are several times their 1.5\,GB Parquet size. OPT-1 gives the same reuse without using the heap.

\subsection{OPT-4: cache only small reused intermediates}
\emph{Bottleneck:} even with OPT-1, the O2 plan scans the trips 4$\times$ and the O3 plan 3$\times$.
\texttt{flows} and the per-hour counts are reused but recomputed for each use. \emph{Hypothesis:} caching them
(330\,k and $\approx$1\,M rows) cuts each operation to one scan. \emph{Change:} \texttt{--opt cache\_agg}.
\emph{Result:} input records O2 188.8$\to$80.8\,M and O3 161.6$\to$80.8\,M. Side effect: the cached data keeps
200 shuffle partitions, which raises O2 to 2983 tasks and its shuffle to 102\,MB. \todo{wall times vs.\ the re-run
of OPT-1+2 under the same load}.

\subsection{OPT-5: shuffle partitions and AQE}
\todo{spark.sql.shuffle.partitions / advisoryPartitionSizeInBytes, to fix the 200 partitions of OPT-4}

\begin{table}[h]
\caption{Summary of optimisations (2019, median of 3 runs).}
\label{tab:summary}
\begin{tabular}{lrrl}
\toprule
configuration & total (s) & vs.\ baseline & note \\
\midrule
baseline & 1398 & -- & \\
+ OPT-1 materialize & 661 & $-$53\,\% & \\
+ OPT-2 agg\_first & 567 & $-$59\,\% & O2 $-$50\,\% \\
OPT-3 cache & fail & -- & OOM, rejected \\
+ OPT-4 cache\_agg & \todo{} & & \\
+ OPT-5 & \todo{} & & \\
\bottomrule
\end{tabular}
\end{table}

%--------------------------------------------------------------------
\section{Analysis}
\todo{Dataset scale: 10\,\%, 20\,\%, 100\,\% of 2019 with the best configuration. Chart: time vs.\ fraction.}
\todo{Cluster size: 1, 2, 3 workers. Speedup, efficiency, why it is not linear.}

%--------------------------------------------------------------------
\section{Discussion}
\todo{Remaining bottlenecks (O3 self-join, which could be a window lag), methodology limits (shared VM, noise,
single machine), what we would do with more time (the full 2016--2020 range).}

%--------------------------------------------------------------------
\newpage
\section*{Acknowledgements and AI usage}
\todo{from docs/ai-usage.md}

\section*{Individual contributions}
\textbf{Tomás Neves (50\,\%):} \todo{Tomás: environment, data, M1, M2, profiler, O2/O3 integration,
optimisation and measurements; confirm the 50/50 split proposed by Hugo.}
\textbf{Hugo Feteira (50\,\%):} M3 (module and step-by-step notebook); O2 and O3 notebooks (design and logic,
later integrated by Tomás); fix of the two O2 bugs found in the audit; synthetic tests and answers to the audit.

\section*{Comments and suggestions}

\bibliographystyle{acm}
\bibliography{refs}
\end{document}
